\chapter{问题与统计图}\label{ch:g5:data}

本章是全年最后一章，综合运用所学知识：多步问题
涉及钱数、数量和时长，以及能一眼看清信息的统计图。
核心思想是“求一”——
即求出单个物品的价格。这一思想将
在\cref{ch:g6:propdata}中发展为比例。

\section{多步问题}

\begin{method}[解决长题目]\label{met:g5:data:steps}
\begin{enumerate}
  \item 读完全题，说出最后要求的问题；
  \item 逆向思考：要回答它，先需要知道什么？
  \item 按顺序逐步求解，每一步写出一个运算和
        一句简短的话；
  \item 将最终答案与常识和估算结果对照，
        检查其是否合理。
\end{enumerate}
\end{method}

\begin{example}\label{ex:g5:data:multistep}
“某校订购 $4$ 箱书，每箱 $25$ 本，每本 $6$ 元，
另加运费 $180$ 元。总费用是多少？”
\begin{enumerate}
  \item 书的数量：$4 \times 25 = 100$ 本。
  \item 书的费用：$100 \times 6 = 600$ 元。
  \item 总费用：$600 + 180 = 780$ 元。
\end{enumerate}
估算检验：约一百本书，每本 $6$ 元，约 $600$ 元，
再加约 $200$ 元，约 $800$ 元——答案 $780$ 元是合理的。
\end{example}

\begin{example}[求一]\label{ex:g5:data:perone}
“$5$ 个相同的杯子共 $12.50$ 元，$8$ 个杯子要多少钱？”
先求出\emph{一}个杯子的价格：
\[
12.50 \div 5 = 2.50,
\qquad\text{然后}\qquad
8 \times 2.50 = 20 .
\]
由“一”出发，一切都可求得——这种两步方法
可解决一类问题（配方、速度、价格），它也是
也是\cref{ch:g6:propdata}的起点。
\end{example}

\begin{example}[找零与预算]\label{ex:g5:data:budget}
佐伊有 $30$ 元。她买了一本书，价 $12.40$ 元，又买了两支笔，
每支 $2.30$ 元。她还能买得起一个 $13$ 元的棋盘游戏吗？
\begin{enumerate}
  \item 笔的费用：$2 \times 2.30 = 4.60$ 元。
  \item 已花费：$12.40 + 4.60 = 17$ 元。
  \item 剩余：$30 - 17 = 13$ 元——正好够买棋盘游戏，
        且没有剩余。
\end{enumerate}
\end{example}

\section{读统计图}

\begin{method}[读任何统计图]\label{met:g5:data:read}
\begin{enumerate}
  \item 读标题：统计的是什么？
  \item 看坐标轴或图例：一格刻度、一个柱子、
        一个符号分别代表什么？
  \item 然后再回答问题——用直尺指着统计图
        有助于准确读出高度。
\end{enumerate}
\end{method}

\begin{example}\label{ex:g5:data:chart}
某镇逐月降雨量：

\begin{omfigure}
\begin{tikzpicture}
\begin{axis}[
  width=8.8cm, height=5.0cm,
  ybar, bar width=13pt,
  symbolic x coords={Jan,Feb,Mar,Apr,May,Jun},
  xtick=data, xticklabels={1月,2月,3月,4月,5月,6月},
  ymin=0, ymax=90,
  ytick={20,40,60,80},
  ylabel={降雨量 (mm)},
  tick label style={font=\small},
  label style={font=\small}]
  \addplot[fill=omDef!30, draw=omDef] coordinates
    {(Jan,60) (Feb,45) (Mar,55) (Apr,70) (May,80) (Jun,30)};
\end{axis}
\end{tikzpicture}

{\small 降雨量柱状图。一格刻度 $= 20$ mm。}
\end{omfigure}

读数：降雨最多的是 5 月（$80$ mm）；6 月的降雨量
是 4 月的\omterm{def:g3:division:half}{一半}吗？（$30$ 对 $70$？不，$70$ 的\omterm{def:g3:division:half}{一半}是 $35$，
所以\emph{并不恰好}是\omterm{def:g3:division:half}{一半}：准确读数很重要！）；六个月的
总降雨量：$60 + 45 + 55 + 70 + 80 + 30 = 340$ mm。
\end{example}

\begin{example}[变化的量]\label{ex:g5:data:line}
一株植物的高度，每周测量一次：

\begin{omfigure}
\begin{tikzpicture}
\begin{axis}[omaxis,
  width=8.4cm, height=5.0cm,
  xmin=0, xmax=6.5, ymin=0, ymax=24,
  xlabel={周}, ylabel={高度 (cm)},
  xtick={1,...,6}, ytick={5,10,15,20}]
  \addplot[omDef, thick, mark=*, mark size=1.6pt] coordinates
    {(1,3) (2,6) (3,10) (4,15) (5,19) (6,21)};
\end{axis}
\end{tikzpicture}

{\small 连接各测量点，可以清楚地看出\emph{增长}情况：
线段越陡，该周植物生长得越快。}
\end{omfigure}

植物在第 $3$ 周至第 $4$ 周之间长得最快（增长 $+5$ 厘米），
末尾则放缓，仅增长 $+2$ 厘米（第 $6$ 周）。
\end{example}

\section{练习}

\begin{exercise}[$\star$]\label{exo:g5:data:1}
分步求解：“某俱乐部租一辆大巴车需 $240$，并购买 $32$ 张博物馆门票，
每张售价 $7$。这次出游的总费用是多少？”
\end{exercise}

\begin{exercise}[$\star$]\label{exo:g5:data:2}
用“归一法”求解：“$3$ 千克苹果售价 $7.50$，
$5$ 千克苹果需要多少钱？”
\end{exercise}

\begin{exercise}[$\star$]\label{exo:g5:data:3}
求解：“六位朋友平均分摊一次野餐的费用 $87$，
每人应付多少？”
\end{exercise}

\begin{exercise}[$\star$]\label{exo:g5:data:4}
萨姆有 $25$，他买了一件 T 恤，售价 $13.90$，又买了一顶帽子，售价 $8.50$。
他还剩多少钱？
\end{exercise}

\begin{exercise}[$\star$]\label{exo:g5:data:5}
根据\cref{ex:g5:data:chart}中的降雨量图：哪些月份
降雨量超过 $50$ 毫米？五月比六月多降了
多少毫米？
\end{exercise}

\begin{exercise}[$\star$]\label{exo:g5:data:6}
根据\cref{ex:g5:data:line}中的植物生长图：第 $2$ 周时植物有多高？
植物哪两周之间恰好增长了
$4$ 厘米？
\end{exercise}

\begin{exercise}[$\star$]\label{exo:g5:data:7}
画出某球队各月进球数的条形统计图：
九月 $4$ 球、十月 $7$ 球、十一月 $3$ 球、十二月 $6$ 球。选取合适的
刻度，并求出总进球数。
\end{exercise}

\begin{exercise}[$\star$]\label{exo:g5:data:8}
一包 $8$ 杯酸奶售价 $3.20$；若单独购买，每杯售价
$0.50$。用归一法求出一包中每杯酸奶的价格，并求出
买一整包时每杯可省多少钱？
\end{exercise}

\begin{exercise}[$\star\star$]\label{exo:g5:data:9}
“某影院周六售出 $148$ 张电影票，每张 $9.50$，周日以相同价格售出 $95$ 张电影票。
周六比周日多收入多少钱？”
请用两种方法求解：一是分别计算两天的收入，
二是利用两天售票张数的\emph{差}来计算——
并检验两种结果是否一致。
\end{exercise}

\begin{exercise}[$\star\star$]\label{exo:g5:data:10}
一份食谱供 $4$ 人食用，需要大米 $300$ 克。艾琳要为
$10$ 人做饭，她需要多少大米？（先求每人的用量——
允许出现小数。）
\end{exercise}

\begin{exercise}[$\star\star$]\label{exo:g5:data:11}
编一道用钱的三步应用题，其最终答案为 $5.50$，
并仿照\cref{met:g5:data:steps}的格式写出完整解答，
并在同学中试做检验。
\end{exercise}
